\clearpage
\fancyhead[L]{\small PALLAS.}
\hypertarget{native-465-pallas}{}
\begin{center}{\large PALLAS.}\end{center}

\begin{center}
\rule{0.55\linewidth}{0.4pt}\\[2pt]
Astronomische Nachrichten Band VI. Nr. 124 Seite 67\,..\,68. 1827 Oktober.\\[2pt]
\rule{0.55\linewidth}{0.4pt}
\end{center}

\begin{center}Pallasbeobachtungen.\end{center}

Hofrat \textsc{Gauss} hat mir die folgenden seiner Pallasbeobachtungen mitgeteilt.

\begin{center}
\begin{tabular}{l|l|l|l}
 & & Rektasz. & $\delta$ \\
\hline
1827 Aug. 31 & $11^{\rm u}16^{\rm m}2^{\rm s}0$ & $328^\circ19'54''6$ & $+5^\circ52'27''4$ \\
\hfill Sept. 2 & $11\;6\;42.2$ & $327\;57\;51.4$ & $5\;28\;38.2$ \\
\hfill 3 & $11\;2\;2.8$ & $327\;46\;57.4$ & $5\;16\;31.8$ \\
\hfill 4 & $10\;57\;24.2$ & $327\;36\;14.9$ & $5\;4\;21.3$ \\
\end{tabular}
\end{center}

Die erste Beobachtung ist weniger zuverlässig, weil der Planet aufgrund des großen Temperaturunterschieds zwischen draußen und drinnen äußerst schwer zu sehen war. Bei den 3 anderen Beobachtungen waren die Fenster schon lange vorher geöffnet worden, und Hofrat \textsc{Gauss} hält diese Beobachtungen für sehr gut.
\begin{flushright}\textsc{Schumacher.}\end{flushright}

